\documentclass[11pt]{article}
\usepackage[a4paper,margin=1in]{geometry}
\usepackage{amsmath,amssymb,amsthm,booktabs,microtype}
\usepackage{cleveref}
\newtheorem{theorem}{Theorem}[section]
\newtheorem{corollary}[theorem]{Corollary}
\newtheorem{proposition}[theorem]{Proposition}
\theoremstyle{remark}
\newtheorem{remark}[theorem]{Remark}
\newcommand{\R}{\mathbb R}
\newcommand{\cP}{\mathcal P}
\newcommand{\cQ}{\mathcal Q}
\newcommand{\Tr}{\operatorname{Tr}}
\title{Moving-Order Duhamel Composition:\\Retained Paths and the Trace-Observation Gap}
\author{Bin Wang}
\date{July 2026}
\begin{document}
\maketitle

\begin{abstract}
RH-324 proves a sharp $O(\sigma)$ physical-to-affine error for one endpoint
leg, but a local bound is not yet a moving-order trace theorem.  We give the
missing abstract composition ledger.  For nonautonomous Markov kernels, a
retained-coordinate path telescope bounds the full path-law error by the seed
error plus the conditional row errors evaluated against the correctly
transported incoming laws.  Thus $O(k)$ phase-matched $O(\sigma)$ legs would
produce $O(k\sigma)=o(kR^{-2k})$ on the natural first-alias clock.  A
two-state example proves that accuracy at one fixed seed is not composable
under phase transport.  For weighted traces we state the exact operator
Duhamel criterion and derive the sharp stability window
$\gamma<1-\log R/\log\lambda=0.3503698834\ldots$.  Finally, a cyclic-shift
Markov family has row and retained-path error $2/n\to0$, zero endpoint error,
and trace gap one.  Hence Markov contraction alone cannot control a growing
trace observation.  The second physical leg, observation norms, parity,
neighboring shell, and joint first-alias trace law remain open.
\end{abstract}

\section{Retained-path Markov composition}

Let $X_0,\ldots,X_m$ be standard Borel spaces.  For $1\le j\le m$, let
$P_j$ and $Q_j$ be Markov kernels from $X_{j-1}$ to $X_j$.  They represent
the physical and affine rows.  Given entrance probabilities $\mu$ and $\nu$
on $X_0$, define the retained-coordinate path laws
\begin{align}
 \cP(dx_0\cdots dx_m)
 &=\mu(dx_0)\prod_{j=1}^mP_j(x_{j-1},dx_j),\label{eq:path-P}\\
 \cQ(dx_0\cdots dx_m)
 &=\nu(dx_0)\prod_{j=1}^mQ_j(x_{j-1},dx_j).\label{eq:path-Q}
\end{align}
All $L^1$ distances below are unhalved; total variation is one half of them.
Put $\nu^P_{j}=\nu P_1\cdots P_j$, with $\nu^P_0=\nu$.

\begin{theorem}[Retained-coordinate Markov Duhamel criterion]
\label{thm:path-duhamel}
The full path laws satisfy
\begin{equation}
 \boxed{
 \|\cP-\cQ\|_1
 \le \|\mu-\nu\|_1
 +\sum_{j=1}^m
 \int_{X_{j-1}}
 \|P_j(x,\cdot)-Q_j(x,\cdot)\|_1\,d\nu^P_{j-1}(x).}
 \label{eq:path-bound}
\end{equation}
Every marginal of the path, including the final $X_m$ law, obeys the same
upper bound.  In particular, if the $j$th row error is uniformly at most
$\varepsilon_j$, then
\begin{equation}
 \|\cP-\cQ\|_1\le\|\mu-\nu\|_1+\sum_{j=1}^m\varepsilon_j.
 \label{eq:sup-bound}
\end{equation}
\end{theorem}

\begin{proof}
First compare the physical path started from $\mu$ with the same physical
path started from $\nu$.  Because $x_0$ is retained and all later factors are
probability kernels, this difference has $L^1$ norm exactly
$\|\mu-\nu\|_1$.

Starting now from $\nu$, telescope through the hybrid path laws whose first
$j$ rows are physical and whose remaining rows are affine.  The $j$th hybrid
difference contains the signed row $P_j-Q_j$.  The prefix has
$x_{j-1}$ marginal $\nu^P_{j-1}$, while the affine suffix has total mass one
for every retained $x_j$.  Tonelli's theorem therefore makes the norm of this
hybrid difference exactly the $j$th integral in \eqref{eq:path-bound}.
The triangle inequality proves the path bound.  Marginalization is an
$L^1$ contraction, and \eqref{eq:sup-bound} follows.
\end{proof}

The retained coordinate is essential: it prevents cancellations between
different source rows before their conditional errors are measured.  The
incoming law is also essential; it is the physical prefix transport of the
reference entrance law, not an arbitrarily reused local seed.

\begin{proposition}[Same-seed accuracy need not compose]
\label{prop:phase-counterexample}
There are two-state Markov kernels $T,P,Q$ and a seed $\delta_0$ such that
\begin{equation}
 \|\delta_0P-\delta_0Q\|_1=0,
 \qquad
 \|\delta_0TP-\delta_0TQ\|_1=2.
 \label{eq:phase-counterexample}
\end{equation}
\end{proposition}

\begin{proof}
Let $T$ send both states to state $1$, let $P$ be the identity, and let $Q$
send both states to state $0$.  At the original seed, both $P$ and $Q$ return
$\delta_0$.  After $T$, the incoming law is $\delta_1$; $P$ returns
$\delta_1$ and $Q$ returns $\delta_0$, whose $L^1$ distance is two.
\end{proof}

This is the exact finite obstruction to iterating a phase-matched estimate
without proving how the phase and profile enter the next row.

\section{The natural-clock consequence}

Use the archived clock and radius
\begin{equation}
 k_\sigma=\frac{\log(1/\sigma)}{2\log\lambda}+O(1),
 \qquad R=1.4,
 \qquad
 \theta=\frac{\log R}{\log\lambda}
 =0.6496301165394707\ldots .
 \label{eq:clock}
\end{equation}
Then $R^{-2k_\sigma}=\Theta(\sigma^\theta)$.

\begin{corollary}[Conditional first-alias negligibility]
\label{cor:markov-clock}
Suppose $m_\sigma=O(k_\sigma)$, the entrance error is
$o(k_\sigma R^{-2k_\sigma})$, and every transported row in
\cref{thm:path-duhamel} has error $O(\sigma)$ with a common constant.  Then
\begin{equation}
 \|\cP_\sigma-\cQ_\sigma\|_1
 =O(k_\sigma\sigma)
 =o(k_\sigma R^{-2k_\sigma}).
 \label{eq:markov-negligible}
\end{equation}
The same holds for every marginal.
\end{corollary}

\begin{proof}
The theorem gives $O(m_\sigma\sigma)=O(k_\sigma\sigma)$.  Dividing by the
target leaves $O(\sigma^{1-\theta})\to0$ because $\theta<1$.
\end{proof}

This corollary is deliberately conditional.  RH-324 proves the required
phase-matched $O(\sigma)$ scale only for the first physical endpoint leg.
It does not prove a common bound for all transported rows or the second
critical physical leg.

\section{Operator products and trace observations}

Markov path control is not yet trace control.  Let $B_0,\ldots,B_m$ be
Banach spaces with $B_m=B_0$, and let $A_j,G_j:B_{j-1}\to B_j$ be bounded
operators.  Suppose a linear observation $\mathfrak t$ on endomorphisms of
$B_0$ obeys
\begin{equation}
 |\mathfrak t(H)|\le T\|H\|.
 \label{eq:observation-bound}
\end{equation}

\begin{theorem}[Trace-observation Duhamel criterion]
\label{thm:operator-duhamel}
The exact identity
\begin{equation}
 A_m\cdots A_1-G_m\cdots G_1
 =\sum_{j=1}^m
 A_m\cdots A_{j+1}(A_j-G_j)G_{j-1}\cdots G_1
 \label{eq:duhamel-identity}
\end{equation}
implies
\begin{equation}
 \left|\mathfrak t(A_m\cdots A_1-G_m\cdots G_1)\right|
 \le\sum_{j=1}^m W_j\delta_j,
 \label{eq:operator-bound}
\end{equation}
where
\begin{align}
 \delta_j&=\|A_j-G_j\|,\nonumber\\
 W_j&=T
 \prod_{\ell=j+1}^m\|A_\ell\|
 \prod_{\ell=1}^{j-1}\|G_\ell\|.
 \label{eq:weights}
\end{align}
Consequently a first-alias trace replacement follows if
$\sum_jW_j\delta_j=o(k_\sigma R^{-2k_\sigma})$.
\end{theorem}

\begin{proof}
Equation \eqref{eq:duhamel-identity} is the finite product telescope.
Apply \eqref{eq:observation-bound}, submultiplicativity, and the triangle
inequality term by term.
\end{proof}

\begin{theorem}[Sharp power window for the Duhamel majorant]
\label{thm:power-window}
Assume $m_\sigma=O(k_\sigma)$, $\delta_j=O(\sigma)$ uniformly, and
\begin{equation}
 \max_jW_j=O(\sigma^{-\gamma}).
 \label{eq:weight-growth}
\end{equation}
Then \eqref{eq:operator-bound} is
$O(k_\sigma\sigma^{1-\gamma})$.  It is
$o(k_\sigma R^{-2k_\sigma})$ whenever
\begin{equation}
 \boxed{
 \gamma<\gamma_*:=1-\theta
 =0.3503698834605293\ldots .}
 \label{eq:gamma-threshold}
\end{equation}
At $\gamma=\gamma_*$ the majorant is only of target order, while for
$\gamma>\gamma_*$ this argument does not decay relative to the target.
\end{theorem}

\begin{proof}
The ratio of the Duhamel majorant to the target has power
\begin{equation}
 \frac{k_\sigma\sigma^{1-\gamma}}
 {k_\sigma\sigma^\theta}
 =\sigma^{1-\theta-\gamma}.
\end{equation}
Its exponent is positive precisely below \eqref{eq:gamma-threshold}.
\end{proof}

RH-18 proves only the lower bound
$\operatorname{cond}(D_k^{\rm G})\ge\sigma^{-1/4+o(1)}$.  The exponent
$1/4$ lies below $\gamma_*$ by
\begin{equation}
 \gamma_*-\frac14=0.1003698834605293\ldots .
 \label{eq:quarter-slack}
\end{equation}
If a relevant stability weight had a matching quarter-power \emph{upper}
bound, this scalar budget would remain compatible.  No such upper bound or
identification with the actual trace observation is archived.

\section{A dimension-free trace bound is false}

Let $S_n$ be the cyclic shift on $n\ge2$ states and put
\begin{equation}
 P_n=I_n,
 \qquad
 Q_n=\left(1-\frac1n\right)I_n+\frac1nS_n.
 \label{eq:trace-example}
\end{equation}
Both matrices are Markov and preserve the uniform law.
Let $\Pi_n^P(i,j)=n^{-1}P_n(i,j)$ and
$\Pi_n^Q(i,j)=n^{-1}Q_n(i,j)$ denote their one-step retained path laws.

\begin{proposition}[Growing-state Markov-to-trace obstruction]
\label{prop:trace-counterexample}
For the pair \eqref{eq:trace-example},
\begin{equation}
 \max_i\|P_n(i,\cdot)-Q_n(i,\cdot)\|_1=\frac2n,
 \qquad
 \|\Pi_n^P-\Pi_n^Q\|_1=\frac2n,
 \label{eq:row-path-small}
\end{equation}
and the endpoint marginals are identical, but
\begin{equation}
 |\Tr P_n-\Tr Q_n|=1.
 \label{eq:trace-gap}
\end{equation}
Thus no dimension-free bound converts Markov row or path $L^1$ error into
trace error.
\end{proposition}

\medskip
\begin{proof}
Each row moves mass $1/n$ from its diagonal entry to the next state, giving
row distance $2/n$.  Averaging the rows against the uniform entrance law
gives the same retained one-step path distance.  Both kernels preserve that
uniform law, so the endpoint error is zero.  Since $S_n$ has zero diagonal,
\[
 \Tr Q_n=n\left(1-\frac1n\right)=n-1,
\]
which proves \eqref{eq:trace-gap}.
\end{proof}

The reproduction table uses $m_\sigma=\lceil2k_\sigma\rceil$, local constant
one, and model growth $\sigma^{-\gamma}$:
\begin{center}\small
\begin{tabular}{c@{\quad}rrrr}
\toprule
$\sigma$ & $\gamma=0$ & $\gamma=1/4$ & $\gamma=\gamma_*$ & $\gamma=0.4$\\
\midrule
$10^{-4}$  &0.080321&0.803212&2.024463&3.197646\\
$10^{-8}$  &0.003187&0.318677&2.024463&5.050691\\
$10^{-12}$ &0.000126&0.126436&2.024463&7.977582\\
\bottomrule
\end{tabular}
\end{center}
The entries are majorant-to-target ratios computed from the exact exponents;
they are not fitted estimates of the physical operator.

\begin{remark}[Claim boundary]
RH-325 proves composition criteria and explicit counterexamples.  It does
not show that every physical leg meets the Markov hypotheses, control the
second critical leg, bound the actual trace observation, combine parity with
the neighboring shell, or prove the joint first-alias matching law.  It
does not prove divergence when the sufficient majorant fails.  No full-trace
replacement is obtained.  Gates A--E remain false/open.  This paper
constructs no Hilbert--Polya operator, identifies no Riemann zero, proves no
von Mangoldt trace formula or zeta-divisor equality, and does not imply RH.
\end{remark}

RH-326 should now derive the parity-renormalized first-alias packet identity
without discarding the transported phase or the sign needed for later joint
cancellation.

\end{document}
